\clearpage\thispagestyle{empty}
\pdfbookmark[0]{Quellen zu Gauß an Bode (21. April 1824)}{bode1824april-sources}
\begin{center}\textbf{Quellen zu Gauß an Bode (21. April 1824)}\end{center}
\noindent Diese KI-Übertragung ins heutige Deutsch folgt Gauß’ Brief vom 21. April 1824 in \href{https://gdz.sub.uni-goettingen.de/id/PPN236007467}{den \textit{Werken}, Band~VI} (Göttingen, 1874), S.~449--450. Bodes \textit{Astronomisches Jahrbuch für 1827} (Berlin, 1824), S.~\href{https://books.google.nl/books?id=4C8PAAAAQAAJ&pg=PA128&redir_esc=y}{128}, \href{https://books.google.nl/books?id=4C8PAAAAQAAJ&pg=PA129&redir_esc=y}{129} und \href{https://books.google.nl/books?id=4C8PAAAAQAAJ&pg=PA130&redir_esc=y}{130}, druckt ihn unter der Überschrift \textit{Meridianbeobachtungen des Kometen von 1824}. Die Links führen zum Digitalisat in Google Books.
\par\smallskip\noindent Der Erstdruck datiert die ersten beiden Beobachtungszeilen auf den 31. Januar und den 3. Februar. In der Werkausgabe stehen der 31. Februar und der 3. Januar. Die Übertragung behält die Datumsangaben der Werkausgabe, die leere Sekundenstelle und den Doppelpunkt in der Zeile für den 20. sowie die beiden Gruppen von je 16 Vergleichszeilen bei. Der Wert $+963.4$ und Gauß’ Vermutung eines Fehlers von einer Zeitminute in der Rektaszension der Altonaer Beobachtung bleiben wie gedruckt erhalten. KI-Übertragung und KI-Satz.
\clearpage
